% from ega4-17-fr.tex
\oldpage[IV-4]{90}

On a vu en outre (16.2.2, (iii)) que $\operatorname{gr}_1(g)$ est ici surjectif, donc on déduit de (17.13.1.2) la suite exacte
\[
\label{IV.17.13.1.3}
\tag{17.13.1.3}
g^*(\sh{N}_{Y/X})\xrightarrow{\alpha}\Omega^1_{X'/S}\otimes_{\sh{O}_{X'}}\sh{O}_{Y'}\to\Omega^1_{Y'/S}\to0.
\]
\end{env}

\begin{proposition}[17.13.2]
\label{IV.17.13.2}
Les notations étant celles de (17.13.1), soit $x'$ un point de $Y'$, $x=g(x')$ son image dans $Y$; supposons $X$ et $Y$ lisses sur $S$ au point $x$, $X'$ lisse sur $S$ au point $x'$. Soient $m$, $m-c$ les dimensions relatives de $X$ et $Y$ sur $S$ au point $x$ (17.10.1), $n$ celle de $X'$ sur $S$ au point $x'$. Les conditions suivantes sont équivalentes:
\begin{enumerate}
  \item[a)] $Y'$ est lisse sur $S$ au point $x'$ et de dimension relative $n-c$.
  \item[b)] L'homomorphisme $\alpha\otimes1:g^*(\sh{N}_{Y/X})\otimes_{\sh{O}_{Y'}}k(x')\to\Omega^1_{X'/S}\otimes_{\sh{O}_{X'}}k(x')$ déduit de l'homomorphisme $\alpha$ de (17.13.1.3) est injectif.
\end{enumerate}
Soit $s$ l'image canonique de $x$ dans $S$, et supposons que $x'$ soit rationnel sur $k(s)$. Alors les conditions a) et b) sont aussi équivalentes à la suivante:
\begin{enumerate}
  \item[b')] L'homomorphisme transposé de $\alpha\otimes1$
  \[
  \mathrm{T}_{X'/S}(x')\to(\sh{N}_{Y/X}\otimes_{\sh{O}_Y}k(x'))^\vee=\mathrm{T}_{X/S}(x)/\mathrm{T}_{Y/S}(x)
  \]
  (cf. (16.5.12)) est surjectif.
\end{enumerate}
De plus, lorsque les conditions a) et b) sont vérifiées en $x'$, elles le sont dans un voisinage de $x'$ dans $Y'$; en remplaçant au besoin $X'$ par un voisinage de $x'$, l'homomorphisme
\[
\operatorname{gr}_1(g):g^*(\sh{N}_{Y/X})\to\sh{N}_{Y'/X'}
\]
est bijectif, et la suite
\[
\label{IV.17.13.2.1}
\tag{17.13.2.1}
0\to g^*(\sh{N}_{Y/X})\xrightarrow{\alpha}\Omega^1_{X'/S}\otimes_{\sh{O}_{X'}}\sh{O}_{Y'}\to\Omega^1_{Y'/S}\to0
\]
obtenue en ajoutant un $0$ à (17.13.1.3), est exacte.
\end{proposition}

Le fait que si les conditions équivalentes a) et b) sont vérifiées au point $x'$, elles le sont aussi au voisinage de $x'$ dans $Y'$, résulte de ce que l'ensemble des points où un morphisme est lisse est ouvert (17.3.7) et de (17.10.2).

En vertu de (17.12.1) appliqué à $X'$ et $Y'$, dire que $Y'$ est lisse sur $S$ au point $x'$ équivaut à dire que l'homomorphisme
\[
\delta\otimes1:\sh{N}_{Y'/X'}\otimes_{\sh{O}_{Y'}}k(x')\to\Omega^1_{X'/S}\otimes_{\sh{O}_{X'}}k(x')
\]
est injectif, compte tenu de (0, 19.1.12) et de ce que $\Omega^1_{X'/S}$ est un $\sh{O}_{X'}$-Module localement libre au point $x'$ (17.2.3); comme alors $\Omega^1_{Y'/S}$ est aussi un $\sh{O}_{Y'}$-Module localement libre au voisinage de $x'$ et que la suite
\[
\label{IV.17.13.2.2}
\tag{17.13.2.2}
0\to\sh{N}_{Y'/X'}\to\Omega^1_{X'/S}\otimes_{\sh{O}_{X'}}\sh{O}_{Y'}\to\Omega^1_{Y'/S}\to0
\]
est exacte (17.2.5), dire que $Y'$ est de dimension relative $n-c$ sur $S$ au point $x'$ signifie, par (17.10.2), que $\sh{N}_{Y'/X'}$ (qui est localement libre au voisinage de $x'$) est de rang $c$ au point $x'$. Mais en vertu des hypothèses sur $X$ et $Y$ au point $x$, et par le même raisonnement, $\sh{N}_{Y/X}$ est localement libre et de rang $c$ au point $x$, donc $g^*(\sh{N}_{Y/X})$ est aussi localement libre et de rang $c$ au point $x'$. Comme l'homomorphisme $\operatorname{gr}_1(g):g^*(\sh{N}_{Y/X})\to\sh{N}_{Y'/X'}$ est surjectif, les conditions précédentes équivalent à dire que cet homomorphisme est

